\subsection{معماری و پیاده‌سازی}
کلاس موتور فشرده‌سازی \lr{UniQLCompressor} در فایل \lr{Codes/framework/compression\_engine/uniql\_compressor.py} مستقر بوده و تحت نام‌های \lr{uniql} و \lr{uniql\_compressor} در رجیستری مرکزی بارگذاری می‌شود. شاخص‌ترین قابلیت‌های طراحی این موتور عبارتند از:
\begin{itemize}
    \item \textbf{حلقه تنظیم مجدد تطبیقی (\lr{Adaptive Bit \& Rank Retries}):} اجرای فرآیند تکرار شونده تا سقف \lr{max\_retries=3} که در صورت بالا بودن خطای تقریب اولیه، به صورت هوشمند و خودکار نرخ بیت یا رتبه را افزایش می‌دهد تا بازنمایی معنایی لایه حفظ گردد.
    \item \textbf{پشتیبانی ساختاری از لایه‌های ترانهاده \lr{Conv1D}:} شناسایی خودکار ویژگی‌های لایه‌های کانولوشن یک‌بعدی در مدل‌های دیکودر (\lr{GPT-2}) و اعمال مناسب ترانهاده قبل و بعد از تجزیه.
    \item \textbf{محافظت از ساختارهای بحرانی شبکه:} مستثنی کردن خودکار لایه‌های پایانی سر طبقه‌بندی، واژگان و تعبیه از فشرده‌سازی به منظور تثبیت احتمالات خروجی.
    \item \textbf{مسیر امن بازگشت به وزن‌های اصلی (\lr{Safe Fallback}):} جایگزینی لایه اصلی دست‌نخورده در صورتی که خطای نهایی پس از تمامی تلاش‌ها فراتر از آستانه اطمینان باشد.
\end{itemize}

\subsection{محدودیت‌ها و تحلیل شکست}
با وجود نتایج تجربی برجسته، محدودیت‌های تحلیلی این روش به شرح زیر است:
\begin{itemize}
    \item \textbf{شبیه‌سازی کوانتیزاسیون به جای فرمت سخت‌افزاری واقعی (\lr{Fake Quantization}):} پیاده‌سازی فعلی عملگر کوانتیزاسیون را با گردکردن مقادیر در تانسورهای ممیز شناور پیاده می‌سازد؛ بنابراین کاهش واقعی فیزیکی مصرف حافظه در سطح ذخیره‌سازی نوع داده به فرمت عدد صحیح ۸ بیتی سخت‌افزاری (\lr{INT8 INT4 Tensor Cores}) نیازمند الحاق هسته‌های محاسباتی سفارشی است.
    \item \textbf{خطای تجمیعی در شبکه‌های دارای مقادیر پرت شدید:} در صورتی که دامنه‌ی مقادیر وزن در لایه‌ای بسیار گسترده باشد، کوانتیزاسیون یکنواخت مقادیر کوچک را با تقریب درشت گرد می‌کند که می‌تواند منجر به فعال‌سازی مکرر حلقه تلاش مجدد و انتخاب رتبه‌های بالاتر گردد.
\end{itemize}
